\documentclass[11pt]{article}

\usepackage[T1]{fontenc}
\usepackage{lmodern}
\usepackage[a4paper,margin=29mm]{geometry}
\usepackage{amsmath,amssymb,amsthm,mathtools}
\usepackage{microtype}
\usepackage{enumitem}
\usepackage{xcolor}
\usepackage[colorlinks=true,linkcolor=blue!45!black,citecolor=blue!45!black,
  urlcolor=blue!55!black]{hyperref}
\setlength{\emergencystretch}{2em}

\newtheorem{theorem}{Theorem}[section]
\newtheorem{proposition}[theorem]{Proposition}
\newtheorem{lemma}[theorem]{Lemma}
\newtheorem{corollary}[theorem]{Corollary}
\newtheorem{definition}[theorem]{Definition}
\newtheorem{remark}[theorem]{Remark}
\newtheorem{question}[theorem]{Question}

\newcommand{\C}{\mathbb C}
\newcommand{\N}{\mathbb N_0}
\newcommand{\dom}{\operatorname{dom}}
\newcommand{\ran}{\operatorname{ran}}
\newcommand{\Graph}{\operatorname{Graph}}
\newcommand{\clos}{\operatorname{cl}}
\newcommand{\cG}{\mathcal G}
\newcommand{\cM}{\mathcal M}
\newcommand{\ip}[2]{\langle #1,#2\rangle}

\title{Two-Sided Singular Defects in Rigged Hilbert Spaces:\\
Classification and Failure of Automatic Quasi Compatibility}
\author{Michael Haufschild}
\date{}

\begin{document}

\maketitle

\begin{abstract}
Quasi Gelfand triples replace the continuous embeddings of a rigged Hilbert space by closed,
densely defined embeddings.  A pre-quasi triple satisfies the pivot pairing identity but not,
a priori, the maximal adjoint-domain condition.  Two published conjectures assert that the
latter condition is automatic, first for Hilbert spaces and then for reflexive Banach spaces.

We disprove both conjectures.  More structurally, in the standard Hilbert anti-dual model we
classify every negative embedding compatible with a fixed positive embedding.  If
$A$ is the positive embedding and $B=(A^*)^{-1}$, then compatible closed negative embeddings are
exactly the restrictions of $B$ whose graphs are closed subspaces of $\Graph B$ with both
coordinate projections dense.  Quasi compatibility holds exactly for the saturated graph
$\Graph B$.

We then identify the sharp obstruction.  Give $\Graph B$ its graph Hilbert norm and let $P$ and
$Q$ be its domain and range coordinate maps.  A proper closed subspace with both projections
dense exists if and only if neither $P$ nor $Q$ is bounded below.  Equivalently, $B$ is both
unbounded and not bounded below.  Whenever this two-sided singularity occurs, a codimension-one
orthogonal hyperplane already produces a defect.  Outside this regime every pre-quasi restriction
is automatically quasi.

An explicit diagonal construction on $\ell^2(\{0,1\}\times\N)$ has alternating ultraviolet and
infrared singular channels.  A single graph hyperplane has dense domain and dense range but is
proper, giving a Hilbert pre-quasi Gelfand triple which is not quasi.  The classification, sharp
criterion, and counterexample have also been formalized unconditionally in Lean~4.
\end{abstract}

\medskip
\noindent\textbf{2020 Mathematics Subject Classification.}
Primary 46C07, 47A06; Secondary 46A20, 47A70.

\noindent\textbf{Keywords.}
Rigged Hilbert space; Gelfand triple; closed operator; linear relation; graph subspace;
adjoint domain; two-sided singularity.

\section{Introduction}

A classical Gelfand triple consists of a dense continuous embedding of a test space into a
Hilbert pivot space and the corresponding anti-dual embedding.  Closed, unbounded embeddings
arise naturally when the two norms are not globally comparable.  This led Cojuhari and Gheondea
to closely embedded Hilbert triplets \cite{CojuhariGheondea2015}, and Skrepek to quasi Gelfand
triples, a Banach-space framework motivated in part by boundary spaces for multidimensional
differential operators and port-Hamiltonian systems \cite{Skrepek2021,Skrepek2025}.

The definition separates a formal pairing condition from a maximality condition.  A
\emph{pre-quasi} triple has closed, densely defined, injective embeddings $\iota_+$ and
$\iota_-$ with dense ranges and satisfies
\[
  \ip{\iota_-g}{\iota_+f}_{0}=\ip{g}{f}_{-,+}.
  \tag{1.1}
\]
It is \emph{quasi} when additionally
\[
  \ran\iota_-=\dom\iota_+^*.
  \tag{1.2}
\]
The pairing identity always gives inclusion in (1.2); the question is whether a proper closed
restriction can retain dense domain and dense range.  Skrepek formulated the following two
conjectures \cite[Conjectures~6.7 and~6.8]{Skrepek2025}:
\begin{enumerate}[label=(\roman*)]
  \item every Hilbert pre-quasi Gelfand triple is quasi;
  \item every pre-quasi Gelfand triple is quasi.
\end{enumerate}
The first was described as plausible despite unsuccessful proof attempts.  A positive answer
would also have supplied an abstract route from weak to strong boundary trace definitions.

The obstacle is not failure of closedness, adjoint existence, or density separately.  It is a
specific infinite-dimensional geometry: a closed operator graph can contain a proper closed
hyperplane whose two coordinate projections are both dense.  Such a hyperplane is invisible in
each ambient Hilbert topology separately, but visible in graph topology.  The phenomenon requires
singularity in both directions.  If one coordinate controls graph norm, dense projection forces
the hyperplane to be the whole graph.

This paper has three main results.

First, Theorem~\ref{thm:classification} gives the exact graph-subspace classification.  For a
closed, densely defined, injective positive embedding $A$ with dense range, its compatible maximal
negative partner is
\[
  B=(A^*)^{-1}.
  \tag{1.3}
\]
Every compatible negative embedding is a restriction of $B$.  Closed pre-quasi partners are
therefore classified by closed subspaces of $\Graph B$ with dense first and second projections,
and quasi compatibility is precisely graph saturation.

Second, Theorem~\ref{thm:sharp} solves the resulting graph problem.  If $P,Q$ are the coordinate
maps on the graph Hilbert space, a proper closed bi-dense subspace exists exactly when both maps
lack a uniform lower estimate.  One vector outside $\ran P^*\cup\ran Q^*$ produces the required
hyperplane.  In operator language this says exactly
\[
  B\text{ unbounded}\qquad\text{and}\qquad
  \inf_{0\ne x\in\dom B}\frac{\|Bx\|}{\|x\|}=0.
  \tag{1.4}
\]
Thus simultaneous ultraviolet and infrared singularity is not merely sufficient; it is the
complete obstruction.  Corollary~\ref{cor:automatic} is the positive survivor: if either side is
regular, automatic quasi compatibility is true.

Third, Theorem~\ref{thm:counterexample} gives a concrete counterexample.  On two copies of
$\ell^2(\N)$, the canonical operator multiplies one channel by $k+1$ and the other by $(k+1)^{-1}$.
The two channels make (1.4) hold.  A graph hyperplane normal to the square-summable vector
$((k+1)^{-1},(k+1)^{-1})_{k\ge0}$ has dense images under both coordinate maps.  Restricting to it
gives a pre-quasi triple which is not quasi.  Since all three spaces are Hilbert spaces, this
refutes both conjectures.

The result does not assert that weak and strong traces differ for any particular PDE.  It shows
that their equality cannot follow from the abstract pre-quasi axioms alone.  Additional PDE
structure may rule out two-sided singularity; recent equality results under Lipschitz hypotheses
are consistent with our criterion \cite{SkrepekPauly2025}.

\subsection*{Relation to prior work}

Closely embedded Hilbert triplets impose the completed compatibility represented here by the full
graph \cite{CojuhariGheondea2015}.  Skrepek proved that a pre-quasi embedding can always be extended
to the canonical quasi partner \cite[Lemma~4.5]{Skrepek2025}; our classification determines every
proper compatible restriction below that extension and decides exactly when one can remain
bi-dense.  Standard closed-operator theory supplies the graph norm, adjoint, and closed-range
ingredients \cite{Kato1995,Schmudgen2012}.  The new point is their conjunction as an exact defect
classification for pre-quasi Gelfand triples, followed by a sharp existence theorem and a
counterexample to the two explicit conjectures.

\section{Pre-quasi triples and graph relations}

All Hilbert spaces are complex.  Inner products are linear in the first argument.  For an operator
$T$ we distinguish its domain $\dom T$, range $\ran T$, and graph
\[
  \Graph T=\{(x,Tx):x\in\dom T\}.
\]
An operator is called densely ranged when $\clos(\ran T)$ is the target space.  A map $T:E\to F$
between Hilbert spaces is \emph{bounded below in inverse form} if
\[
  \|x\|_E\le C\|Tx\|_F\qquad(x\in E)
  \tag{2.1}
\]
for some $C\ge0$.  For injective maps this is the usual uniform lower bound with reciprocal
constant.

We need only a standard Hilbert anti-dual pair, which is already a subclass of the published weak
conjecture.

\begin{definition}[standard Hilbert pre-quasi pair]\label{def:prequasi}
Let $H$ be a Hilbert space.  A pair of linear operators
\[
  A:\dom A\subset H\to H,
  \qquad C:\dom C\subset H\to H
\]
is a \emph{standard Hilbert pre-quasi pair} if both operators are closed, densely defined,
injective, and densely ranged, and
\[
  \ip{Cg}{Af}=\ip{g}{f}
  \qquad(g\in\dom C,\ f\in\dom A).
  \tag{2.2}
\]
It is \emph{quasi compatible} if $\ran C=\dom A^*$.
\end{definition}

Taking $\mathcal X_+=\mathcal X_0=\mathcal X_-=H$ with its standard complete anti-dual pairing
turns Definition~\ref{def:prequasi} into a pre-quasi Gelfand triple in the sense of
\cite[Definition~4.1]{Skrepek2025}.  Thus a counterexample in this standard model is a
counterexample to the Hilbert conjecture, without any identification or renorming argument.

Suppose from now on that $A$ is closed, densely defined, injective, and densely ranged.  Dense
range gives $\ker A^*=\{0\}$, so the inverse relation of $A^*$ is single-valued.  We write
\[
  B=(A^*)^{-1},\qquad
  \dom B=\ran A^*,\qquad
  \ran B=\dom A^*.
  \tag{2.3}
\]

\begin{proposition}[canonical partner]\label{prop:canonical}
The operator $B$ in \emph{(2.3)} is closed, densely defined, injective, and densely ranged.  It
satisfies
\[
  \ip{Bg}{Af}=\ip{g}{f}
  \qquad(g\in\dom B,\ f\in\dom A).
  \tag{2.4}
\]
If $C$ satisfies the pairing identity \emph{(2.2)}, then $C\subset B$ as operators.
\end{proposition}

\begin{proof}
Because $A$ is densely defined and closed, $A^*$ is closed and densely defined.  Since
$\clos(\ran A)=H$, the identity $(\ran A)^\perp=\ker A^*$ gives $\ker A^*=\{0\}$; hence its inverse
$B$ is an operator and is closed.  The domain of $B$ is $\ran A^*$, which is dense because
$(\ran A^*)^\perp=\ker A=\{0\}$.  Its range is $\dom A^*$, which is dense.  The inverse is injective
because $A^*$ is single-valued.

For $g=A^*y\in\dom B$ one has $Bg=y$, and therefore
\[
  \ip{Bg}{Af}=\ip{y}{Af}=\ip{A^*y}{f}=\ip{g}{f}.
\]
Now let $g\in\dom C$.  Equation (2.2) says that the functional
$f\mapsto\ip{Cg}{Af}$ on $\dom A$ is represented by $g$.  By the definition and uniqueness of the
adjoint, $Cg\in\dom A^*$ and $A^*Cg=g$.  Thus $g\in\ran A^*=\dom B$ and $Bg=Cg$, proving $C\subset B$.
\end{proof}

The graph of $B$ is a Hilbert space
\[
  \cG_B=(\Graph B,\|\cdot\|_B),\qquad
  \|(x,Bx)\|_B^2=\|x\|^2+\|Bx\|^2.
  \tag{2.5}
\]
Let
\[
  P:\cG_B\to H,\quad P(x,Bx)=x,
  \qquad
  Q:\cG_B\to H,\quad Q(x,Bx)=Bx.
  \tag{2.6}
\]
Both are bounded, injective, and densely ranged.  Their ranges are respectively $\dom B$ and
$\ran B$.

\begin{definition}
A closed subspace $M\subset\cG_B$ is \emph{bi-dense} when $P(M)$ and $Q(M)$ are dense in $H$.
It is a \emph{graph defect} when in addition $M\ne\cG_B$.
\end{definition}

This use of ``defect'' concerns a proper subrelation of a single-valued closed operator.  It is
distinct from the multivalued part of a general linear relation.

\section{Classification by closed bi-dense graph subspaces}

\begin{theorem}[graph-subspace classification]\label{thm:classification}
Fix $A$ as above and put $B=(A^*)^{-1}$.  The assignments
\[
  C\longmapsto \Graph C,
  \qquad
  M\longmapsto C_M,\quad
  \Graph C_M=M,
  \tag{3.1}
\]
give inverse, inclusion-preserving correspondences between
\begin{enumerate}[label=(\alph*)]
  \item closed operators $C$ for which $(A,C)$ is a standard Hilbert pre-quasi pair;
  \item closed bi-dense subspaces $M$ of $\cG_B$.
\end{enumerate}
Moreover,
\[
  (A,C)\text{ is quasi compatible}
  \quad\Longleftrightarrow\quad
  C=B
  \quad\Longleftrightarrow\quad
  \Graph C=\cG_B.
  \tag{3.2}
\]
\end{theorem}

\begin{proof}
Let $(A,C)$ be pre-quasi.  Proposition~\ref{prop:canonical} gives $C\subset B$, hence
$\Graph C\subset\cG_B$.  Closedness of $C$ makes $\Graph C$ closed in $H\oplus H$ and therefore
closed in $\cG_B$.  Its first projection is $\dom C$ and its second projection is $\ran C$; both
are dense.  Thus $\Graph C$ is bi-dense.

Conversely, let $M\subset\cG_B$ be closed and bi-dense.  Since the first coordinate map is
injective on an operator graph, $M$ is the graph of an operator $C_M$.  It is closed in
$H\oplus H$.  Its domain and range are $P(M)$ and $Q(M)$, hence dense.  Because $C_M\subset B$ and
$B$ is injective, $C_M$ is injective.  Restricting (2.4) gives the pairing identity (2.2), so
$(A,C_M)$ is pre-quasi.  The two constructions in (3.1) are inverse and preserve inclusion.

It remains to identify saturation.  Proposition~\ref{prop:canonical} gives
$\ran B=\dom A^*$.  Hence $C=B$ implies quasi compatibility.  Conversely, suppose
$C\subset B$ and $\ran C=\dom A^*=\ran B$.  For $x\in\dom B$, choose $d\in\dom C$ with
$Cd=Bx$.  Since $C\subset B$ and $B$ is injective, $Bd=Bx$ implies $d=x$.  Thus
$\dom B\subset\dom C$, so $C=B$.
\end{proof}

\begin{remark}
The canonical extension $C\mapsto B$ is known in the pre-quasi framework
\cite[Lemma~4.5]{Skrepek2025}.  Theorem~\ref{thm:classification} identifies the entire interval
below that extension: maximality fails exactly through a closed bi-dense subrelation.  The issue
is therefore not whether a quasi completion exists, but whether density in two weaker coordinate
topologies forces graph-topology saturation.
\end{remark}

\section{The sharp two-sided singularity criterion}

We first isolate two Hilbert-space lemmas.  Their combination converts graph defects into a
precise lower-bound obstruction.

\begin{lemma}[adjoint-surjectivity criterion]\label{lem:adjoint-surj}
Let $T:G\to H$ be a bounded injective operator between Hilbert spaces with dense range.  Then
\[
  T^*\text{ is surjective}
  \quad\Longleftrightarrow\quad
  \exists C\ge0\ \forall x\in G:\ \|x\|_G\le C\|Tx\|_H.
  \tag{4.1}
\]
\end{lemma}

\begin{proof}
Suppose $T^*$ is onto.  By the open mapping theorem, there is $C\ge0$ such that every $z\in G$
has a preimage $y\in H$ with $T^*y=z$ and $\|y\|\le C\|z\|$.  Then
\[
  \|z\|^2=|\ip{z}{T^*y}|=|\ip{Tz}{y}|
  \le C\|Tz\|\|z\|,
\]
and cancellation proves the estimate when $z\ne0$.

Conversely, the estimate makes $\ran T$ closed.  Since it is dense, $T$ is surjective.  Thus $T$
is a bounded isomorphism and $T^{-1}$ is bounded.  From
$T^*(T^{-1})^*=I_G$ we obtain surjectivity of $T^*$.
\end{proof}

\begin{lemma}[dense hyperplane image]\label{lem:hyperplane}
Let $T:G\to H$ be bounded with dense range, and let $v\in G\setminus\ran T^*$.  Then
\[
  T(v^\perp)\text{ is dense in }H.
  \tag{4.2}
\]
\end{lemma}

\begin{proof}
Take $y\perp T(v^\perp)$.  Then $T^*y\in(v^\perp)^\perp=\C v$, so
$T^*y=\alpha v$.  If $\alpha\ne0$, then $v\in\ran T^*$, contrary to assumption.  Hence
$T^*y=0$.  Dense range of $T$ makes $T^*$ injective, so $y=0$.  Therefore
$T(v^\perp)^\perp=\{0\}$, which is equivalent to (4.2).
\end{proof}

\begin{definition}[two-sided singularity]\label{def:twosided}
Let $P:G\to H_1$ and $Q:G\to H_2$ be bounded, injective, densely ranged Hilbert-space operators.
The pair $(P,Q)$ is \emph{two-sided singular} if neither map satisfies an estimate of the form
\emph{(2.1)}.
\end{definition}

\begin{theorem}[sharp defect criterion]\label{thm:sharp}
Let $P:G\to H_1$ and $Q:G\to H_2$ be bounded, injective, and densely ranged.  The following are
equivalent.
\begin{enumerate}[label=(\roman*)]
  \item $(P,Q)$ is two-sided singular.
  \item There exists a proper closed subspace $M\subset G$ such that $P(M)$ is dense in $H_1$ and
        $Q(M)$ is dense in $H_2$.
  \item Such an $M$ can be chosen as a codimension-one orthogonal hyperplane.
\end{enumerate}
\end{theorem}

\begin{proof}
Assume (i).  Lemma~\ref{lem:adjoint-surj} says that $P^*$ and $Q^*$ are not onto.  Choose
$x\notin\ran P^*$ and $y\notin\ran Q^*$.  There is a vector outside both adjoint ranges.  Indeed,
if $x\notin\ran Q^*$, take $v=x$; if $y\notin\ran P^*$, take $v=y$; in the remaining case
$x\in\ran Q^*$ and $y\in\ran P^*$, and $v=x+y$ belongs to neither range.  For example, if
$x+y\in\ran P^*$ then subtracting $y\in\ran P^*$ would put $x$ in $\ran P^*$, a contradiction;
the $Q^*$ argument is symmetric.  Such $v$ is nonzero because zero belongs to both ranges.

Set $M=v^\perp$.  It is a proper closed hyperplane.  Applying Lemma~\ref{lem:hyperplane} to $P$
and $Q$ shows that both $P(M)$ and $Q(M)$ are dense.  This proves (iii), hence (ii).

Conversely, suppose (ii) and assume, for contradiction, that $P$ satisfies
$\|z\|_G\le C\|Pz\|$.  The restriction $P|_M$ has closed range because $M$ is complete and the
same lower estimate holds on $M$.  Its range is also dense, hence $P(M)=H_1$.  Given $z\in G$,
choose $m\in M$ with $Pm=Pz$.  Injectivity of $P$ gives $m=z$, so $M=G$, contrary to properness.
Thus $P$ is not bounded below in inverse form.  The same argument applies to $Q$, proving (i).
\end{proof}

Apply Theorem~\ref{thm:sharp} to the coordinate maps (2.6).

\begin{corollary}[operator form]\label{cor:operator}
Let $B$ be closed, densely defined, injective, and densely ranged.  A graph defect exists if and
only if
\[
  B\text{ is unbounded and }B\text{ is not bounded below}.
  \tag{4.3}
\]
Whenever it exists, a codimension-one graph defect exists.
\end{corollary}

\begin{proof}
For $z=(x,Bx)\in\cG_B$, inverse boundedness of $P$ is the estimate
\[
  (\|x\|^2+\|Bx\|^2)^{1/2}\le C\|x\|.
\]
It holds exactly when $B$ is bounded on its domain.  Since $B$ is closed and densely defined,
this is equivalent to $\dom B=H$ and $B\in\mathcal B(H)$: a bounded estimate extends $B$ to $H$,
and closedness identifies the extension with $B$.

Inverse boundedness of $Q$ is
\[
  (\|x\|^2+\|Bx\|^2)^{1/2}\le C\|Bx\|,
\]
which is equivalent to $\|x\|\le C'\|Bx\|$, namely that $B$ is bounded below.  Theorem
\ref{thm:sharp} now gives the claim.
\end{proof}

\begin{corollary}[maximal positive survivor]\label{cor:automatic}
Fix a standard Hilbert pre-quasi pair $(A,C)$ and let $B=(A^*)^{-1}$.  If $B$ is bounded or bounded
below, then $(A,C)$ is quasi compatible.  Equivalently, automatic quasi compatibility fails for a
fixed positive embedding $A$ precisely when its canonical partner $B$ is two-sided singular.
\end{corollary}

\begin{proof}
Theorem~\ref{thm:classification} turns a non-quasi pair into a graph defect.  Corollary
\ref{cor:operator} excludes such a defect if either regularity condition holds.  Conversely, if
$B$ is two-sided singular, Theorem~\ref{thm:sharp} produces a proper closed bi-dense hyperplane;
Theorem~\ref{thm:classification} turns it into a non-quasi pre-quasi partner.
\end{proof}

The corollary is sharp for every canonical graph, not merely for the diagonal example below.

\section{A diagonal counterexample}

We now realize the two singular directions on disjoint channels.  This makes every domain,
range, and properness assertion explicit.

Let
\[
  H=\ell^2(\{0,1\}\times\N),\qquad r_k=k+1.
\]
For $m=(m_{0,k},m_{1,k})$ define bounded diagonal contractions $P,Q\in\mathcal B(H)$ by
\[
\begin{array}{lll}
  (Pm)_{0,k}=r_k^{-1}m_{0,k},&& (Pm)_{1,k}=m_{1,k},\\[1mm]
  (Qm)_{0,k}=m_{0,k},&& (Qm)_{1,k}=r_k^{-1}m_{1,k}.
\end{array}
\tag{5.1}
\]
They are positive, self-adjoint, injective, and densely ranged.  Define $B$ by
\[
  \dom B=\ran P,\qquad B(Pm)=Qm.
  \tag{5.2}
\]
Then $B$ is the maximal diagonal operator
\[
  (Bx)_{0,k}=r_kx_{0,k},\qquad
  (Bx)_{1,k}=r_k^{-1}x_{1,k},
  \tag{5.3}
\]
on the exact maximal domain
\[
  \dom B=\left\{x\in H:
    \sum_{k\ge0}\big(r_k^2|x_{0,k}|^2+r_k^{-2}|x_{1,k}|^2\big)<\infty\right\},
\]
and its inverse $A=B^{-1}$ is the maximal diagonal operator
\[
  \dom A=\ran Q,\qquad A(Qm)=Pm,
  \tag{5.4}
\]
with the two diagonal weights interchanged and
\[
  \dom A=\left\{x\in H:
    \sum_{k\ge0}\big(r_k^{-2}|x_{0,k}|^2+r_k^2|x_{1,k}|^2\big)<\infty\right\}.
\]
Both are positive self-adjoint operators by the maximal real diagonal multiplication theorem.

For completeness, closedness can be seen without invoking spectral theory.  The graph
parameterization
\[
  R:H\to H\oplus H,\qquad Rm=(Pm,Qm)
  \tag{5.5}
\]
has the bounded left inverse
\[
  S(x,y)=\big((y_{0,k})_{k\ge0},(x_{1,k})_{k\ge0}\big).
  \tag{5.6}
\]
The range of $R$ is the intersection of the closed coordinate constraints
$x_{0,k}=r_k^{-1}y_{0,k}$ and $y_{1,k}=r_k^{-1}x_{1,k}$, hence it is closed and equals
$\Graph B$.  The finite-support vectors show density of all relevant domains and ranges.

Now put
\[
  v_{0,k}=v_{1,k}=r_k^{-1},\qquad M=v^\perp\subset H.
  \tag{5.7}
\]
The vector $v$ belongs to $H$, since $\sum_{k\ge0}r_k^{-2}<\infty$.  It belongs to neither
$\ran P$ nor $\ran Q$: a $P$-preimage would have constant $0$-channel equal to one, while a
$Q$-preimage would have constant $1$-channel equal to one.  Neither constant channel is in
$\ell^2$.

Define the restriction
\[
  \dom C=P(M),\qquad C(Pm)=Qm\quad(m\in M).
  \tag{5.8}
\]

\begin{theorem}[explicit counterexample]\label{thm:counterexample}
With $A$ and $C$ defined by \emph{(5.4)} and \emph{(5.8)}, the pair $(A,C)$ is a standard Hilbert
pre-quasi pair but is not quasi compatible.  Consequently Conjectures~6.7 and~6.8 of
\cite{Skrepek2025} are false.
\end{theorem}

\begin{proof}
The graph of $C$ is $R(M)$.  It is closed: if $Rm_n\to z$ in $H\oplus H$, then
$m_n=S(Rm_n)\to S(z)$; closedness of $M$ gives $S(z)\in M$, and continuity of $R$ yields
$z=R(S(z))\in R(M)$.  Thus $C$ is closed.  It is injective because $Q$ is injective.

The operators $P$ and $Q$ are self-adjoint.  Since $v\notin\ran P=\ran P^*$ and
$v\notin\ran Q=\ran Q^*$, Lemma~\ref{lem:hyperplane} gives
\[
  \clos(P(M))=H,\qquad \clos(Q(M))=H.
  \tag{5.9}
\]
Hence $C$ has dense domain and dense range.  The operator $A$ is self-adjoint, closed, injective,
dense-domain, and dense-range by its diagonal description.

For $g=Pm\in\dom C$ with $m\in M$ and $f=Qn\in\dom A$ with $n\in H$, the commuting positive
diagonals $P,Q$ give
\[
  \ip{Cg}{Af}=\ip{Qm}{Pn}=\ip{Pm}{Qn}=\ip{g}{f}.
  \tag{5.10}
\]
Thus $(A,C)$ is pre-quasi.

Because $A=A^*$, one has $\dom A^*=\dom A=\ran Q$.  On the other hand,
$\ran C=Q(M)$.  The hyperplane $M$ is proper and $Q$ is injective, so $Q(M)$ is a proper subset of
$\ran Q$.  Therefore $\ran C\ne\dom A^*$ and the pair is not quasi compatible.

Finally, $H$ paired with itself is a complete Hilbert dual pair, so $(H,H,H)$ with embeddings
$A,C$ satisfies every hypothesis of the weak Hilbert conjecture.  Its failure refutes that
conjecture.  The strong conjecture includes the Hilbert case and is therefore false as well.
\end{proof}

\begin{remark}[even--odd form]
Any bijection $\{0,1\}\times\N\simeq\N$ transports the construction to ordinary
$\ell^2(\N)$.  Under the even--odd bijection, $B$ has diagonal
$1,1,2,\frac12,3,\frac13,\ldots$ up to the harmless ordering of the first two entries.  The split
index makes the two singular channels and the bounded recovery map (5.6) transparent.
\end{remark}

\section{Consequences and boundaries}

\subsection{Exact logical status of the conjectures}

Theorem~\ref{thm:counterexample} is a counterexample inside the smallest convenient category:
all three spaces coincide as Hilbert spaces, the duality is the standard Hilbert duality, and the
canonical positive and negative operators are positive self-adjoint diagonal operators.  Failure
therefore does not depend on exotic Banach duality, non-reflexivity, a multivalued adjoint, or a
nonclosable map.

The defect is also minimal in graph codimension.  Deleting a one-dimensional graph-normal
direction is enough, while each coordinate image stays dense.  Theorem~\ref{thm:sharp} shows that
this is not an accident of diagonal weights: every two-sided singular canonical graph has a
codimension-one defect, and no other graph can have any proper closed bi-dense defect.

\subsection{What survives positively}

The failed conjectures contain a sharp true theorem.  If the canonical inverse adjoint is bounded,
the domain coordinate controls graph norm.  If it is bounded below, the range coordinate controls
graph norm.  Either estimate upgrades density of a closed compatible restriction to graph
saturation.  Only simultaneous failure permits nonmaximality.

For ordinary Gelfand triples, one embedding is continuous and one is obtained by completion; this
places the construction on a one-sided regular branch.  The present theorem explains why the
classical framework does not display the defect and why merely replacing continuity by closedness
opens it.

\subsection{Boundary traces}

Skrepek observed that a positive solution of the weak conjecture would imply equality of weak and
strong boundary trace definitions in the motivating differential-operator setting
\cite{Skrepek2025}.  The logical converse does not follow.  Our result proves that pre-quasi
axioms alone cannot establish that equality.  A concrete PDE may carry coercivity, a Poincar\'e
estimate, closed-range structure, or a regular decomposition that makes one coordinate bounded
below.  Such an estimate falls exactly under Corollary~\ref{cor:automatic}.  Accordingly,
PDE-specific weak-equals-strong theorems remain untouched \cite{SkrepekPauly2025}.

\subsection{Relation-theoretic interpretation}

The canonical inverse adjoint $B=(A^*)^{-1}$ is most naturally viewed through its closed graph
relation.  Pairing compatibility places every admissible negative graph inside this relation.
The two coordinate topologies see only $P(M)$ and $Q(M)$, whereas closedness sees $M$ in graph
norm.  Two-sided singularity is exactly non-equivalence of graph norm with either coordinate norm.
Thus the obstruction is a topology mismatch within a single-valued relation, not a multivalued
singular part produced by graph closure.

This interpretation suggests a useful diagnostic for applications: identify the canonical graph,
test its two coordinate lower estimates, and only then ask whether a dense compatible core is
maximal.  Density in both pivot images is decisive only after one coordinate controls graph norm.

\section{Formal verification}

The complete result has been formalized in Lean~4 against the repository's pinned Mathlib
toolchain.  The formal development uses closed partially defined linear maps and their adjoints;
no finite-dimensional surrogate, numerical evaluation, unproved density assumption, or extra
maximality hypothesis is introduced.

The principal checked declarations are:
\begin{itemize}\small\raggedright
  \item \texttt{preQuasi\_graph\_le\_canonical\_graph}, and\\
        \texttt{preQuasi\_of\_closed\_biDense\_}\\
        \texttt{canonicalGraphSubspace}, the two directions of
        Theorem~\ref{thm:classification};
  \item \texttt{preQuasi\_quasi\_iff\_}\\
        \texttt{eq\_canonicalNegative}, the saturation equivalence (3.2);
  \item \texttt{twoSidedSingular\_iff\_}\\
        \texttt{hasProperClosedBiDenseSubspace}, the sharp equivalence in
        Theorem~\ref{thm:sharp};
  \item \texttt{exists\_hilbertPreQuasiPair\_not\_quasi}, the unconditional counterexample of
        Theorem~\ref{thm:counterexample}.
\end{itemize}
The explicit formal model uses $\ell^2(\mathrm{Bool}\times\N)$, which is the channel notation of
Section~5.  Axiom inspection of all five headline declarations reports only Lean's standard
logical axioms \texttt{propext}, \texttt{Classical.choice}, and \texttt{Quot.sound}.

\section{Further questions}

The obstruction theorem closes automaticity under the published axioms, but it opens more
structured questions.

\begin{question}
Which analytic hypotheses on boundary maps force one coordinate projection of the canonical graph
to be bounded below?  Corollary~\ref{cor:automatic} reduces any positive PDE theorem of this type
to a coercive or closed-range estimate.
\end{question}

\begin{question}
For two-sided singular canonical graphs, which additional invariances restrict the lattice of
closed bi-dense subspaces?  The present hyperplane construction is universal, but locality,
differential order, or symmetry may exclude its use as an admissible operator core.
\end{question}

\begin{question}
Can the graph-subspace classification be extended from standard Hilbert anti-duality to reflexive
Banach dual pairs with an equally sharp geometric criterion?  The Hilbert proof uses orthogonal
hyperplanes and adjoint ranges at its decisive step, so a Banach analogue should involve norming
annihilators and closed-range geometry rather than a formal transcription.
\end{question}

\section*{Acknowledgment of scope}

The paper resolves the automaticity conjectures negatively and classifies the exact Hilbert graph
mechanism.  It does not claim a counterexample to weak--strong trace equality for a specified
differential operator, nor a Banach-space classification beyond the Hilbert standard-pair model.

\begingroup
\raggedright
\begin{thebibliography}{99}

\bibitem{CojuhariGheondea2015}
P.~Cojuhari and A.~Gheondea,
\emph{Triplets of closely embedded Hilbert spaces},
Integral Equations Operator Theory \textbf{81} (2015), 1--33.
\href{https://doi.org/10.1007/s00020-014-2195-0}{doi:10.1007/s00020-014-2195-0}.

\bibitem{Kato1995}
T.~Kato,
\emph{Perturbation Theory for Linear Operators},
Classics in Mathematics, Springer, Berlin, 1995.
\href{https://doi.org/10.1007/978-3-642-66282-9}{doi:10.1007/978-3-642-66282-9}.

\bibitem{Schmudgen2012}
K.~Schm\"udgen,
\emph{Unbounded Self-adjoint Operators on Hilbert Space},
Graduate Texts in Mathematics 265, Springer, Dordrecht, 2012.
\href{https://doi.org/10.1007/978-94-007-4753-1}{doi:10.1007/978-94-007-4753-1}.

\bibitem{Skrepek2021}
N.~Skrepek,
\emph{Well-posedness of linear first order port-Hamiltonian systems on multidimensional spatial
domains},
Evolution Equations and Control Theory \textbf{10} (2021), 965--1006.
\href{https://doi.org/10.3934/eect.2020098}{doi:10.3934/eect.2020098}.

\bibitem{Skrepek2025}
N.~Skrepek,
\emph{Quasi Gelfand triples},
Integral Equations Operator Theory \textbf{97} (2025), article 1.
\href{https://doi.org/10.1007/s00020-024-02780-9}{doi:10.1007/s00020-024-02780-9}.

\bibitem{SkrepekPauly2025}
N.~Skrepek and D.~Pauly,
\emph{Weak equals strong $L^2$ regularity for partial tangential traces on Lipschitz domains},
J. Math. Anal. Appl. \textbf{550} (2025), 129548.
\href{https://doi.org/10.1016/j.jmaa.2025.129548}{doi:10.1016/j.jmaa.2025.129548}.

\end{thebibliography}
\endgroup

\end{document}
